% GENERATED FILE. Edit canonical lessons, then run npm run generate:books.
% canonical-source-hash: fnv1a64:764c9af2fce51e45
% canonical-lessons: IT-W260-my-own-line

\chapter{Writing --- My Own Line}
\label{ch:it-my-own-line}
\hlchaptermodality{260}

\begin{chapteropening}
\textbf{By the end of this chapter:} \emph{I can write one line about myself from a closed supply of known words, with nothing to copy, and put the accent on every è that means is.}

Complete IT-W260-my-own-line: a name and one more fact, chosen rather than copied, with every e checked for is or and before anything else.
\end{chapteropening}

\section[(my own line)]{Your own line — the first writing with no model}

\label{lesson:IT-W260-my-own-line}



\begin{quote}
Everything you have written so far had something to reach for. You traced \textbf{ciao}, copied it, wrote it with the model covered, and wrote it from the sound.

This does not. You choose the words, from a supply that is closed: only what this book has taught you.
\end{quote}

\subsection*{Writing — controlled composition}
Write \textbf{one line} about yourself, on paper: your name, then one more fact about you. You hold two ways to give a name — the reflexive verb, and \emph{my name is} with its article kept. Choose one; nobody is choosing for you.

One thing to decide deliberately rather than by accident before you start: \textbf{the accent on è is a letter, not decoration.}

With its accent, \textbf{è} is \emph{is}. Without it, \textbf{e} is \emph{and}. A hand that drops the accent has not written Italian untidily: \emph{il mio nome e Marco} says \textquotedblleft{}my name and Marco\textquotedblright{}, and leaves the reader to guess the verb you meant.

Write the line. Then read it back and check every \textbf{e} — does it mean \emph{is} or \emph{and}? — \textbf{before} you check anything else.

\begin{tcolorbox}[breakable,colback=teal!4,colframe=teal!35!black,title={Before you move on}]
Is the accent on \textbf{è} decoration? (\textbf{No} — without it the word is \emph{and}.)

And what did you choose here that the earlier lessons chose for you? (\textbf{The words.})
\end{tcolorbox}
